
\section{Introduction}\label{sec:intro}

Despite the promising performance of deep neural networks achieved in various practical tasks~\citep{simonyan2014very, ren2015faster, hinton2012deep, mikolov2010recurrent, alipanahi2015predicting, litjens2017survey}, uncertainty quantification (UQ) has attracted growing attention in recent years to fulfill the emerging demand for more robust and reliable machine learning models, as UQ aims to measure the reliability of the model's prediction quantitatively. Accurate uncertainty estimation is especially critical for the field that is highly sensitive to error prediction, such as autonomous driving~\citep{bojarski2016end} and medical diagnosis~\citep{begoli2019need}.

Most state-of-the-art approaches~\citep{gal2016dropout, lakshminarayanan2017simple, malinin2018predictive, van2020simple} focus on building a deep model equipped with uncertainty quantification capability so that a single deep model can achieve both desired prediction and UQ performance simultaneously. However, such an approach for UQ suffers from practical limitations because it either requires a specific model structure or explicitly training the entire model from scratch to impose the uncertainty quantification ability.  A more realistic scenario is to quantify the uncertainty of a pretrained model in a post-hoc manner due to practical constraints. For example, (1) compared with prediction accuracy and generalization performance, uncertainty quantification ability of deep learning models are usually considered with lower priority, especially for profit-oriented applications, such as recommendation systems; (2) some applications require the models to impose other constraints, such as fairness or privacy, which might sacrifice the UQ performance; (3) for some applications such as transfer learning, the pretrained models are usually available, and it might be a waste of resources to train a new model from scratch.  

Motivated by these practical concerns, we focus on tackling the post-hoc uncertainty learning problem, i.e., given a pretrained model, how to improve its UQ quality without affecting its predictive performance. %Compared to the traditional setting, post-hoc uncertainty learning has two significant advantages. First, it is computationally efficient, i.e., it does not require retraining the whole neural network. Second, it is more effective, i.e., when there exists a trade-off between prediction accuracy and uncertainty calibration, post-hoc uncertainty learning avoids this conflict by decomposing the learning process into two sub-tasks, one focusing on prediction performance, the other aiming for UQ.
%The post-hoc uncertainty learning problem is still under-explored, and few works focus on tackling this problem.
Prior works on improving uncertainty quality in a post-hoc setting have been mainly targeted towards improving calibration \citep{guo2017calibration,kull2019beyond}. These approaches typically fail to augment the pre-trained model with the ability to capture different sources of uncertainty, such as epistemic uncertainty, which is crucial for applications such Out-of-Distribution (OOD) detection. Several recent works \citep{chen2019confidence,jain2021deup} have adopted the meta-modeling approach, where a meta-model is trained to predict whether or not the pretrained model is correct on the validation samples. These methods still rely on a point estimate of the meta-model parameters, which can be unreliable, especially when the validation set is small. 
% Furthermore, they tend to provide a confidence measure at the overall prediction level which may not desirable for example in a multi-class setting where we expect the accurate confidence scores for each possible class label.   %Meta-model-based approaches train an auxiliary model to oversee the performance of the pretrained base model. White-box~\citep{chen2019confidence} trains multiple linear classifiers on intermediate layers of the pretrained model using additional validation data, then feed the outputs into another regression model to output a confidence score for predicting whether or not the pretrained model is correct on the validation samples. 
%Similarly, DEUP~\citep{jain2021deup} trains an error predictor to estimate the epistemic uncertainty of the pretrained model in terms of the difference between generalization error and aleatoric uncertainty. 
%Bayesian approaches leverage approximation inference techniques without affecting the pretrained model prediction accuracy, such as Laplace approximation. LULA~\citep{kristiadi2021learnable} train additional hidden units building on layers of MAP-trained model to improve its uncertainty calibration with Laplace approximation. However, all these methods require additional data samples that are either validation or out of distribution dataset to train or tune the hyper-parameter, which is infeasible when these data are not available. 
%Moreover, they are not flexible enough to quantify and distinguish different types of uncertainties, such as epistemic uncertainty and aleatoric uncertainty, which are known to be crucial in various learning applications~\citep{hullermeier2021aleatoric}. 

In this paper, we propose a novel Bayesian meta-model-based uncertainty learning approach to mitigate the aforementioned issues. Our proposed method requires no additional data other than the training dataset and is flexible enough to quantify different kinds of uncertainties and easily adapt to different application settings.
%\textbf{Summary of contributions:} 
%First, we formulate the post-hoc uncertainty learning problem with multiple essential components, which has the flexibility to adapt to different problem settings, such as transfer learning. Next, we propose a simple and effective uncertainty learning method based on three
Our empirical results provides crucial insights regarding meta-model training: (1) The diversity in feature representations across different layers is essential for uncertainty quantification, especially for out-of-domain (OOD) data detection tasks; (2) Leveraging the Dirichlet meta-model to capture different uncertainties, including total uncertainty and epistemic uncertainty; (3) There exists an over-fitting issue in uncertainty learning similar to supervised learning that needs to be addressed by a novel validation strategy to achieve better performance. Furthermore, we show that our proposed approach has the flexibility to adapt to various applications, including OOD detection, misclassification detection, and trustworthy transfer learning. %Finally, we demonstrate our proposed meta-model approach's leading uncertainty quantification performance on multiple representative image classification benchmarks.


% First, we formulate the post-hoc uncertainty learning problem into a general framework with multiple essential components with the flexibility to adapt to different problem settings, such as transfer learning. Next, we propose a simple and effective uncertainty learning method based on three crucial insights: (1) The use of multiple layers of feature representation is important for uncertainty quantification, especially for OOD tasks. (2) Leveraging Dirichlet techniques to model different kinds of uncertainties, such as epistemic uncertainties. (3) There exists an over-fitting issue in uncertainty learning similar to supervised learning that needs to be addressed to achieve desired performance. Then, we show that our proposed approach has large flexibility to adapt to various applications, including OOD detection, misclassification detection, and Trustworthy transfer learning. Finally, we demonstrate the leading uncertainty quantification performance of our proposed meta-model approach on multiple representative image classification benchmarks.
